\documentclass[10pt,a4paper]{amsart}
\usepackage[margin=19mm]{geometry}
\usepackage{amsmath,amssymb,mathtools}
\usepackage[scr=rsfs,cal=euler]{mathalfa}
\usepackage{xfrac}
\usepackage[unicode,hidelinks,bookmarks=false]{hyperref}
\newcommand{\ie}{i.e.\ }
\newcommand{\cf}{cf.\ }
\newcommand{\resp}{respectively\ }
\newcommand{\Subsection}{\subsection}
\providecommand{\nobreakdash}{\nobreak\hbox{-}}
\setlength{\emergencystretch}{2em}
\multlinegap0pt
\usepackage{bm}
\usepackage{bbm}
\usepackage{dsfont}
\usepackage{xspace}
\usepackage{cancel}
\usepackage{mathtools}
\usepackage{amsthm}
\makeatletter
\@ifundefined{theorem}{\newtheorem{theorem}{Theorem}}{}
\@ifundefined{definition}{\newtheorem{definition}[theorem]{Definition}}{}
\@ifundefined{lemma}{\newtheorem{lemma}[theorem]{Lemma}}{}
\makeatother
\title[On cutoff via rigidity for high dimensional curved diffusion]{On cutoff via rigidity for high dimensional curved diffusions}
\author{Djalil Chafa\"i, Max Fathi}
\date{Source: arXiv:2412.15969v4 (CC BY 4.0)}
\begin{document}
\maketitle

\noindent\emph{Source and licence.} On cutoff via rigidity for high dimensional curved diffusions. Version arXiv:2412.15969v4 (CC BY 4.0), distributed under the Creative Commons Attribution 4.0 International licence, as recorded in the arXiv licence metadata for this exact version and shown on the abstract page https://arxiv.org/abs/2412.15969v4. Full rights notice: licenses/arxiv-2412.15969-CC-BY-4.0.txt.\par\smallskip

\noindent\textit{Editorial scope.} Counted unit: the Lemma whose source label is \texttt{lem\_proj\_ou} in the article (source0.tex lines 905--911, source label \texttt{lem\_proj\_ou}), delivered together with the article's own complete original proof of it (source0.tex lines 913--932, one \texttt{proof} environment of 20 lines). No mathematics is written, repaired or re-derived: statement and proof are the article's own text line for line. Required context retained uncounted: eq:L (lines 102--104); df:multeig (lines 667--675); thm\_split\_eucl (lines 866--879). Every definition, display and label of those lines is retained unchanged. Omitted from this package and not redistributed: 1--101, 105--666, 676--865, 880--904, 912--912, 933--1699. Nothing inside a retained excerpt is omitted or altered: every retained body is delivered exactly as the article's lines are written, so no omission receipt of any kind accompanies this package. Citations: the retained excerpts carry no citation command at all, so no bibliography entry of the article is transcribed and no reference is redistributed. Reference adaptation: none was needed and none was made; every reference inside the delivered unit resolves inside it, so no unresolved reference or citation is produced. Printed slips kept literal: no printed word, symbol, hypothesis or number of the counted statement or of its proof was altered. Layout and class: the article's own class is \texttt{amsart} with packages including geometry, hyperref, fourier, amssymb, amsfonts, mathrsfs, amsmath, amsthm, enumitem, natbib; the wrapper is the lane's pinned \texttt{amsart} editorial wrapper at 10pt on A4, which restates the theorem-like environments the retained text needs and adds the article's own macro definitions that the retained text needs (none), with extra wrapper packages (bm, bbm, dsfont, xspace, cancel, mathtools). The delivered numbering is the unit's own. Assets: no retained span contains a figure environment, a graphics-inclusion command, a PDF-page inclusion, a table or a TikZ picture; no image file is copied into this package, the package declares no asset dependency and no image file is redistributed with it. Licence: version arXiv:2412.15969v4 of this article is distributed under the Creative Commons Attribution 4.0 International licence, as recorded in the arXiv licence metadata for this exact version and shown on the abstract page https://arxiv.org/abs/2412.15969v4. A full rights notice is delivered at licenses/arxiv-2412.15969-CC-BY-4.0.txt. Characters: every retained excerpt is pure ASCII, so the delivered engine, XeLaTeX with its default Latin Modern text font, reports no missing character.
\begin{equation}\label{eq:L}
\mathcal{L} = \Delta - \nabla V \cdot \nabla
\end{equation}

\begin{definition}[Multi-eigenfunction]\label{df:multeig}
  Let $E_1$ be the eigenspace of $-\mathcal{L}$ associated with the eigenvalue
  $\lambda_1$, in $L^2(\mu)$. Let $k_1 = \dim{E_1}$. We define a
  multi-eigenfunction as being a map $M \to\mathbb{R}^{k_1}$ whose coordinates are
  orthogonal elements of $E_1$ in $L^2(\mu)$. We denote by $F_1$ the set of
  all multi-eigenfunctions, and by $\mathbb{S}_{F_1}$ the set of multi-eigenfunctions
  whose coordinates are elements of the unit sphere of $E_1$ in $L^2(\mu)$,
  thus orthonormal.
\end{definition}

\begin{theorem}[Gaussian factorization in the Euclidean
  space] \label{thm_split_eucl} Consider the operator \eqref{eq:L} and assume
  that for some constant $\rho>0$, $V$ is $\rho$-convex and
  $\lambda_1 = \rho$. Then there is an orthonormal basis $(e_1,\ldots,e_d)$ of
  $\mathbb{R}^d$ and a vector $m \in \mathbb{R}^d$ such that $V$ is of the form
    \begin{equation}
    V(x) = \frac{\rho}{2}((x_1-m_1)^2 + \cdots + (x_{k_1}-m_{k_1})^2) + \widetilde{V}(x_{k_1+1},\ldots,x_d)
  \end{equation}
  where $k_1$ is the dimension of the eigenspace $E_1$ of $-\mathcal{L}$
  associated with $\lambda_1$, and $\widetilde{V}$ is $\rho$-convex on
  $\mathbb{R}^{d-k_1}$. Moreover, all eigenfunctions with eigenvalue
  $\lambda_1$ are affine, and only depend on the first $k_1$ coordinates in
  the above basis.
\end{theorem}

\begin{lemma}[OU process associated to an eigenfunction] \label{lem_proj_ou}
  Under the setting of Theorem~\ref{thm_split_eucl} and
  Definition~\ref{df:multeig}, for all $T \in \mathbb{S}_{F_1}$, the process
  ${(T(X_t))}_{t\geq0}$ is a $k_1$-dimensional OU process scaled by a factor
  $\rho$, that is a process on $\mathbb{R}^{k_1}$ with generator
  $\rho \Delta - \rho x \cdot \nabla$.
\end{lemma}

\begin{proof}
  Let us show that $T(X)$ is a Markov process and compute its
  generator. %Let us compute the generator.
  If $\vec{v}$ is a vector-valued function, then $\vec{\mathcal{L}} \vec{v} $
  is the vector obtained by applying $\mathcal{L}$ to each coordinate, and
  $\Gamma(\vec{v})$ the matrix whose coefficients are
  $\Gamma(\vec{v}_i, \vec{v}_j)$. Since the coordinates of $T$ are orthogonal
  normalized eigenfunctions, $\vec{\mathcal{L}} T = - \lambda_1 T$ and
  $\Gamma(T) = \rho \operatorname{Id}$, as per Theorem~\ref{thm_split_eucl}
  (affine with orthogonal constant gradients of norm $\sqrt{\rho}$). By the
  diffusion property, we have
  \begin{align}
    \mathcal{L}(g \circ T) &= \nabla g \circ T \cdot \vec{\mathcal{L}} T + \langle \nabla^2 g \circ T, \Gamma(T)\rangle \\
                  &= - \lambda_1 \nabla g\circ T \cdot T + \rho \Delta g \circ T \\
                  &= - \rho \nabla g\circ T \cdot T + \rho \Delta g \circ T.
  \end{align}
  This is a function of $T$, so $T(X)$ is a Markov process, and when viewing
  it as such it is indeed equal to $\rho$ times the generator of an OU process
  with unit variance, applied to a function $g$.
\end{proof}

\end{document}
